\documentclass[12pt,a4paper]{extarticle}
\usepackage[utf8]{inputenc}
\usepackage[T5]{fontenc}
\usepackage{geometry}
\geometry{a4paper, left=1.5cm, right=1cm, top=1.5cm, bottom=1.5cm, headheight=20pt, headsep=15pt, footskip=28pt}
\usepackage{enumitem}
\usepackage{tikz}
\usetikzlibrary{patterns, patterns.meta} % Thêm patterns.meta vào đây
% Gọi package nằm trong thư mục con template
\usepackage{../template/mngocbox}

\begin{document}

% Sử dụng ngoặc vuông [...] để thay đổi linh hoạt các thông số
\makeheadbox[headerright={Tổng ôn 10-11},chude={CHỦ ĐỀ 02},tieude={Bất phương trình bậc nhất hai ẩn, hệ bất phương trình bậc nhất hai ẩn và Quy hoạch tuyến tính},dinhdang={Đại số},lop={12 },gv={Nguyễn Lê Minh Ngọc}]

\vspace{2em}
\makeheading{III}{Đề kiểm tra}
\textbf{Câu 1.} Bạn Khải làm một bài thi giữa học kỳ I môn Toán. Đề thi gồm $35$ câu hỏi trắc nghiệm và $3$ câu hỏi tự luận. Khi làm đúng mỗi câu trắc nghiệm được $0{,}2$ điểm, làm đúng mỗi câu tự luận được $1$ điểm. Giả sử bạn Khải làm đúng $x$ câu trắc nghiệm và $y$ câu tự luận và số điểm bạn Khải đạt được là ít nhất $9$ điểm. Các mệnh đề sau đúng hay sai?
    \begin{enumerate}
        \item Bạn Khải được $0{,}2.x$ điểm trắc nghiệm và $y$ điểm tự luận ($x, y > 0$).
        \item Bất phương trình bậc nhất hai ẩn là $0{,}2x + y \ge 9$.
        \item Cặp $(8; 2)$ là một nghiệm của bất phương trình bậc nhất cho hai ẩn $x, y$.
        \item Bạn Khải trả lời được ít nhất $25$ câu trắc nghiệm và $2$ câu tự luận.
    \end{enumerate}
\textbf{Câu 2.} Một cửa hàng bán hai loại gạo, loại I bán mỗi kg lãi $3000$ đồng, loại II bán mỗi kg lãi $2000$ đồng. Giả sử cửa hàng bán $x$ kg gạo loại I và $y$ kg gạo loại II. Bất phương trình biểu thị mối liên hệ giữa $x$ và $y$ để cửa hàng đó thu được số lãi lớn hơn $100000$ đồng có dạng $ax + by > 10$. Khi đó $a + b$ bằng?\\
\vspace{0.5cm}
\textbf{Câu 3.} Một gian hàng trưng bày bàn và ghế rộng $60\text{ m}^2$. Diện tích để kê một chiếc ghế là $0{,}5\text{ m}^2$, một chiếc bàn là $1{,}2\text{ m}^2$. Gọi $x$ là số chiếc ghế, $y$ là số chiếc bàn được kê.
    \begin{enumerate}
        \item Viết bất phương trình bậc nhất hai ẩn $x, y$ cho phần mặt sàn để kê bàn và ghế, biết diện tích mặt sàn dành cho lưu thông tối thiểu là $12\text{ m}^2$.
        \item Chỉ ra ba nghiệm của bất phương trình trên.
    \end{enumerate}
\textbf{Câu 4.} Hà, Châu, Liên và Ngân cùng đi mua trà sữa. Cả bốn bạn có tất cả $185$ nghìn đồng. Bốn bạn mua $4$ cốc trà sữa với giá tiền $35$ nghìn đồng một cốc. Các bạn gọi thêm trân châu cho vào trà sữa. Một phần trân châu đen có giá $5$ nghìn đồng, một phần trân châu trắng có giá $10$ nghìn đồng.
\textbf{Câu 5.} Một cửa hàng bán lẻ bán hai loại cà phê. Loại thứ nhất 140 nghìn đồng/kg, loại thứ hai 180 nghin đồng/kg. Cửa hàng trộn x kg loại thứ nhất và y kg loại thứ hai sao cho hạt cà phê đã trộn cso giá không quá 170 nghìn đồng/kg
\begin{itemize}
    \item Viết bất phương trình bậc nhất hai ẩn x, y thoả mãn điều kiện đề bài.
    \item Biểu diễn miền nghiệm của bất phương trình tìm được ở câu a trên mặt phẳng toạ độ.
\end{itemize}
\end{document}